\chapter*{Resumen del Trabajo Terminal}
\thispagestyle{fancy}

El presente Trabajo Terminal describe el diseño, análisis y desarrollo parcial de un sistema integral para la gestión de finanzas personales, compuesto por una aplicación móvil multiplataforma y una plataforma web. El sistema permite a los usuarios registrar sus ingresos y gastos de forma automatizada mediante el uso de Reconocimiento Óptico de Caracteres (OCR), el cual extrae la información contenida en tickets de compra físicos o digitales, eliminando la necesidad de ingreso manual de datos.

La solución emplea una arquitectura cliente-servidor en tres capas, con un backend desarrollado en Spring Boot, un frontend web en Next.js y una aplicación móvil construida con Flutter. La información procesada se almacena en una base de datos relacional MySQL y se intercambia mediante una API REST segura. Para la categorización automática de gastos se aplican reglas heurísticas basadas en el modelo de sobres de Dave Ramsey, y los resultados se presentan al usuario a través de tableros interactivos con gráficos dinámicos.

Durante el Trabajo Terminal I se completaron las fases de análisis de requerimientos, diseño del sistema, implementación del frontend web 
y primeras pruebas del módulo OCR. Las fases de implementación del backend, desarrollo de la aplicación móvil, integración de componentes
y pruebas exhaustivas se realizarón en el Trabajo Terminal II.

El proyecto atiende una problemática identificada en la población mexicana: según la Encuesta Nacional sobre Salud Financiera (ENSAFI) 2023, solo el 53.2\,\% de los ciudadanos lleva algún tipo de registro de sus gastos. La solución propuesta busca fomentar hábitos financieros saludables mediante tecnología accesible y de fácil uso, orientada principalmente a jóvenes y adultos sin conocimientos contables avanzados.

\newpage
